\section{Ví dụ thực tiễn}
\subsection{Ví dụ thực tiễn về quy luật giá trị thặng dư trong các quốc gia tư bản chủ nghĩa}

\subsubsection{Khái quát sự vận hành của quy luật giá trị thặng dư trong nền kinh tế tư bản chủ nghĩa}

Trong lý luận của C. Mác, quy luật giá trị thặng dư được xem là quy luật kinh tế cơ bản của phương thức sản xuất tư bản chủ nghĩa. Quy luật này phản ánh mục tiêu tối cao của nền sản xuất tư bản chủ nghĩa là tối đa hóa giá trị thặng dư thông qua việc khai thác sức lao động của người lao động \cite{gt}; \cite{Capital}.

Trong nền kinh tế tư bản chủ nghĩa, người lao động không sở hữu tư liệu sản xuất nên buộc phải bán sức lao động của mình cho nhà tư bản để nhận tiền lương. Tuy nhiên, giá trị mà người lao động tạo ra trong quá trình sản xuất thường lớn hơn giá trị sức lao động mà họ nhận được dưới dạng tiền lương. Phần chênh lệch giữa giá trị mới do lao động tạo ra và tiền lương trả cho người lao động chính là giá trị thặng dư, và phần giá trị này bị nhà tư bản chiếm đoạt dưới dạng lợi nhuận \cite{Capital}.

Giá trị thặng dư không chỉ là nguồn gốc của lợi nhuận mà còn là cơ sở của quá trình tích lũy tư bản và mở rộng sản xuất trong nền kinh tế tư bản chủ nghĩa. Vì vậy, các doanh nghiệp tư bản luôn tìm cách gia tăng giá trị thặng dư thông qua nhiều phương thức khác nhau như kéo dài thời gian lao động, tăng cường độ lao động, cải tiến công nghệ hoặc mở rộng thị trường \cite{gt}.

Trong thực tiễn lịch sử, quy luật giá trị thặng dư thể hiện rõ rệt trong quá trình phát triển của nhiều quốc gia tư bản. Hai trường hợp tiêu biểu thường được nhắc đến là sự phát triển của chủ nghĩa tư bản ở Pháp trong thời kỳ công nghiệp và thuộc địa, và sự hình thành chủ nghĩa tư bản độc quyền – đế quốc ở Hoa Kỳ vào cuối thế kỷ XIX và thế kỷ XX.

\subsubsection{Sự vận động của quy luật giá trị thặng dư trong lịch sử phát triển của chủ nghĩa tư bản}

\textbf{Chủ nghĩa tư bản công nghiệp và thời kỳ thuộc địa ở Pháp}

Trong thế kỷ XIX, cùng với sự lan rộng của Cách mạng Công nghiệp ở châu Âu, nền kinh tế tư bản chủ nghĩa tại Pháp phát triển mạnh mẽ. Các ngành công nghiệp như dệt may, luyện kim, khai khoáng và sản xuất máy móc được mở rộng nhanh chóng. Sự phát triển của hệ thống đường sắt và các trung tâm công nghiệp lớn đã tạo điều kiện cho quá trình tích lũy tư bản diễn ra với tốc độ cao \cite{lichsu}.

Trong giai đoạn này, lực lượng lao động công nghiệp gia tăng nhanh chóng khi nhiều nông dân rời bỏ nông thôn để làm việc trong các nhà máy tại thành thị. Điều kiện lao động của công nhân thường rất khắc nghiệt: thời gian lao động có thể kéo dài từ 10 đến 14 giờ mỗi ngày, trong khi tiền lương chỉ đủ đáp ứng những nhu cầu sinh hoạt tối thiểu \cite{lichsu}.

Theo phân tích của Karl Marx, đây là giai đoạn mà các nhà tư bản chủ yếu tạo ra giá trị thặng dư bằng phương thức giá trị thặng dư tuyệt đối, tức là kéo dài thời gian lao động hoặc tăng cường độ lao động của công nhân trong khi tiền lương không thay đổi tương ứng \cite{gt}; \cite{Capital}.

Bên cạnh sự phát triển của công nghiệp trong nước, Pháp còn mở rộng hệ thống thuộc địa rộng lớn tại nhiều khu vực như Bắc Phi, Tây Phi và đặc biệt là Đông Dương thuộc Pháp. Các thuộc địa này trở thành nguồn cung cấp nguyên liệu, lao động giá rẻ và thị trường tiêu thụ cho các doanh nghiệp Pháp \cite{lichsu}.

Tại Đông Dương, nhiều hoạt động kinh tế quy mô lớn đã được tổ chức như khai thác than ở Quảng Ninh, trồng cao su ở Nam Kỳ và khai thác nhiều loại tài nguyên khác. Lực lượng lao động bản địa thường phải làm việc trong điều kiện khó khăn với mức tiền công thấp. Phần lớn giá trị do lao động tại các thuộc địa tạo ra được chuyển về chính quốc dưới dạng lợi nhuận cho các công ty và nhà tư bản Pháp \cite{lichsu}.

Như vậy, trong thời kỳ này, quá trình tạo ra giá trị thặng dư không chỉ diễn ra trong phạm vi lãnh thổ quốc gia mà còn mở rộng ra hệ thống thuộc địa, qua đó giúp các nước tư bản như Pháp tích lũy một lượng lớn tư bản phục vụ cho quá trình công nghiệp hóa.

\textbf{Chủ nghĩa tư bản độc quyền và chủ nghĩa đế quốc ở Hoa Kỳ}

Vào cuối thế kỷ XIX và đầu thế kỷ XX, chủ nghĩa tư bản thế giới bước sang giai đoạn phát triển mới – giai đoạn tư bản độc quyền. Theo phân tích của Vladimir Lenin trong tác phẩm \textit{Imperialism: The Highest Stage of Capitalism}, giai đoạn này được đặc trưng bởi sự tập trung tư bản cao độ, sự hình thành các tập đoàn độc quyền và xu hướng xuất khẩu tư bản ra nước ngoài \cite{LeninImperial}.

Hoa Kỳ là một trong những quốc gia tiêu biểu cho giai đoạn phát triển này. Từ cuối thế kỷ XIX, nhiều tập đoàn công nghiệp khổng lồ đã xuất hiện, chẳng hạn như Standard Oil trong ngành dầu mỏ, U.S. Steel trong ngành luyện kim và Ford Motor Company trong ngành sản xuất ô tô.

Các tập đoàn này có quy mô vốn rất lớn và kiểm soát phần lớn thị trường trong các ngành công nghiệp quan trọng. Nhờ áp dụng sản xuất quy mô lớn, dây chuyền công nghiệp và tổ chức lao động khoa học, các doanh nghiệp có thể tăng mạnh năng suất lao động, từ đó tạo ra lượng giá trị thặng dư ngày càng lớn \cite{LeninImperial}.

Một ví dụ nổi bật là mô hình sản xuất hàng loạt của Henry Ford trong ngành công nghiệp ô tô đầu thế kỷ XX. Việc áp dụng dây chuyền sản xuất đã làm giảm đáng kể thời gian sản xuất mỗi chiếc xe, đồng thời tăng năng suất lao động của công nhân. Theo quan điểm của Marx, đây chính là hình thức tạo ra giá trị thặng dư tương đối, thông qua việc tăng năng suất lao động và giảm giá trị sức lao động trong tổng thời gian lao động \cite{gt}; \cite{Capital}.

Ngoài ra, trong thế kỷ XX, các tập đoàn Mỹ cũng mở rộng hoạt động ra nhiều quốc gia khác thông qua đầu tư trực tiếp và kiểm soát thị trường toàn cầu. Điều này giúp các doanh nghiệp Mỹ thu được lợi nhuận lớn từ nhiều khu vực khác nhau trên thế giới, qua đó củng cố vị thế kinh tế của Hoa Kỳ trong hệ thống tư bản chủ nghĩa quốc tế.

\textbf{Sự biến đổi của hình thức tạo ra giá trị thặng dư trong chủ nghĩa tư bản hiện đại}

Bước sang nửa sau thế kỷ XX, đặc biệt là sau Chiến tranh thế giới thứ hai, nền kinh tế của nhiều quốc gia tư bản phát triển như Pháp và Hoa Kỳ tiếp tục chuyển biến mạnh mẽ dưới tác động của tiến bộ khoa học – kỹ thuật và quá trình toàn cầu hóa kinh tế. Trong giai đoạn này, mặc dù hình thức biểu hiện có nhiều thay đổi, quy luật giá trị thặng dư vẫn tiếp tục giữ vai trò trung tâm trong quá trình vận động của nền kinh tế tư bản chủ nghĩa \cite{gt}.

Một đặc điểm nổi bật của giai đoạn này là sự phát triển nhanh chóng của khoa học – công nghệ và việc áp dụng rộng rãi tự động hóa, tin học hóa và các phương pháp quản lý hiện đại vào sản xuất. Điều này giúp các doanh nghiệp nâng cao năng suất lao động, giảm chi phí sản xuất và tăng khả năng tạo ra giá trị thặng dư. Theo lý luận của Marx, đây chính là biểu hiện của giá trị thặng dư tương đối, tức là gia tăng phần thời gian lao động thặng dư thông qua việc nâng cao năng suất lao động xã hội \cite{Capital}.

Tại Hoa Kỳ, trong nửa sau thế kỷ XX, nhiều tập đoàn lớn trong các ngành công nghiệp, tài chính và công nghệ tiếp tục mở rộng quy mô sản xuất và ảnh hưởng ra phạm vi toàn cầu. Các doanh nghiệp này không chỉ khai thác lao động trong nước mà còn mở rộng đầu tư ra nhiều quốc gia khác nhằm tận dụng nguồn lao động rẻ, tài nguyên và thị trường mới. Điều này giúp các tập đoàn tư bản tiếp tục thu được lợi nhuận lớn và duy trì quá trình tích lũy tư bản ở quy mô toàn cầu \cite{LeninImperial}.

Tại Pháp, nền kinh tế cũng trải qua quá trình hiện đại hóa mạnh mẽ với sự phát triển của các ngành công nghiệp công nghệ cao, dịch vụ tài chính và các tập đoàn đa quốc gia. Sự phát triển của các doanh nghiệp lớn và quá trình quốc tế hóa sản xuất cho thấy việc tạo ra giá trị thặng dư ngày càng gắn liền với sự phát triển của khoa học – công nghệ và sự mở rộng của thị trường thế giới \cite{lichsu}.

Như vậy, có thể thấy rằng mặc dù điều kiện kinh tế – xã hội và trình độ phát triển khoa học – kỹ thuật đã thay đổi đáng kể so với giai đoạn đầu của chủ nghĩa tư bản, mục tiêu cơ bản của sản xuất tư bản chủ nghĩa vẫn không thay đổi, đó là tối đa hóa giá trị thặng dư thông qua việc tổ chức sản xuất ngày càng hiệu quả và mở rộng phạm vi tích lũy tư bản.

\subsubsection{Nhận xét chung}

Qua hai trường hợp tiêu biểu là Pháp và Hoa Kỳ, có thể thấy quy luật giá trị thặng dư luôn giữ vai trò trung tâm trong sự vận động của nền kinh tế tư bản chủ nghĩa.

Trong giai đoạn tư bản công nghiệp ở Pháp, giá trị thặng dư chủ yếu được tạo ra thông qua việc kéo dài thời gian lao động và khai thác lao động trong các nhà máy cũng như trong hệ thống thuộc địa. Đến giai đoạn tư bản độc quyền ở Hoa Kỳ, việc áp dụng công nghệ, sản xuất quy mô lớn và tổ chức lao động khoa học đã giúp các doanh nghiệp tạo ra giá trị thặng dư chủ yếu thông qua tăng năng suất lao động.

Điều đó cho thấy mặc dù hình thức biểu hiện có thể thay đổi theo từng giai đoạn lịch sử, mục tiêu cơ bản của sản xuất tư bản chủ nghĩa vẫn là tối đa hóa giá trị thặng dư và mở rộng quá trình tích lũy tư bản.